\section{Discussion}
\label{sec:discussion}

\textbf{A unified account.}~Our results address a longstanding question: how energy-efficient computation can coexist with highly variable neuronal firing. The Gibbs-manifold formulation shows that sparseness and variability are coupled descriptors of the same response state via $S = \mathrm{CV}^2/(1+\mathrm{CV}^2)$, while efficiency is organized along the manifold by the marginal efficiency $\beta$.

Optimal codes lie along the one-dimensional Gibbs manifold in the $H$--$C$ plane \cite{laughlin_energy_2001, Harris2012, Shoval2012}; the three systems' trajectories align with this structure (Fig.~5), with ICE above $\beta$ throughout. V1 spontaneous activity occupies the steepest segment, consistent with the irregular resting state maintaining high coding readiness within the prevailing energy budget \cite{Raichle2001}.

These findings complement recent evidence that energy constraints impose low-dimensional structure at the level of intrinsic neuronal parameters \cite{Sommer2026} and classical information-rate analyses showing sublinear gains with firing rate \cite{Rieke1997,Sterling2015}. Our framework instead asks how a fixed budget constrains the entire response \emph{distribution}, suggesting the sublinear information-rate relationship may reflect operation near the Gibbs manifold.

\paragraph{Relation to prior work.}
Four bodies of related work deserve comparison.  \emph{Unified efficient coding theory} \cite{Chalk2018} integrates efficient, predictive, and sparse coding within one information-theoretic optimization; we instead isolate the metabolic-budget axis and characterize the full response-distribution shape under a single Gibbs optimization with CV as the macroscopic coordinate.  \emph{Asynchronous irregular (AI) state theory} \cite{Renart2010,LitwinKumar2012} shows that balanced E--I produces high-CV Poisson-like spiking in recurrent networks; we add a normative rationale, identifying the AI state as energetically efficient and quantifying its efficiency along the Gibbs manifold.  \emph{Maximum-entropy models} of neural populations \cite{45} share our variational structure but fit observed joint distributions under pairwise-correlation constraints, whereas we derive a normative cost-constrained Gibbs manifold with CV as a scalar coordinate.  \emph{Neural-manifold} approaches identify low-dimensional subspaces of population activity; the Gibbs manifold is instead a one-dimensional theoretical curve in the scalar $H$--$C$ plane---the two are complementary, describing stimulus-driven covariance vs.\ marginal response-statistics efficiency, respectively.

\textbf{Thermodynamic reading and the free-energy principle.}~The Gibbs form invites a thermodynamic reading---structural, not dynamical: neural circuits are not thermal systems and we make no claim of detailed balance or microscopic reversibility.  The variational structure (entropy maximization under an expected-cost constraint) is mathematically equivalent to constrained free-energy minimization $F = C - \beta^{-1}H$, with $F$ a neural analogue of Helmholtz free energy.  The Lagrange multiplier $\beta$ is formally an inverse temperature, $dF/dp = 0$ recovers the Gibbs distribution, and the $\beta$-parameterized manifold maps onto an equilibrium thermodynamic state space in this formal sense alone.

This connection resonates with but differs from the free energy principle (FEP) of Friston \cite{friston_free-energy_2010}: FEP minimizes \emph{variational} free energy to maintain homeostasis and model the external world, whereas we minimize a \emph{thermodynamic} free energy to maximize information per unit cost.  Both select exponential-family distributions but constrain different quantities (divergence from a generative model vs.\ metabolic expenditure); the two may converge where metabolic efficiency and accurate world-modeling are jointly optimized---an intersection meriting further investigation.

\textbf{From biological mechanisms to engineering principles.}~The geometry of the Gibbs manifold is fixed by the optimization problem, independent of implementation.  Synaptic plasticity and E-I balance determine \emph{where} the system operates along it, not the manifold's shape.  Our models show that local plasticity rules---inhibitory STDP \cite{16,sp_stdp} and Hebbian/anti-Hebbian learning---suffice to steer circuits toward the high-CV, high-efficiency regime \cite{43}, consistent with theories linking E-I balance to criticality \cite{Liang2025prl} but adding a thermodynamic dimension.

The same principle translates into an engineering rule.  Replacing hand-tuned sparsity priors with CV-maximization recovers V1-like receptive fields at a more favourable efficiency--fidelity Pareto position than the Cauchy-prior baseline under matched conditions, because the RE-gated CV objective avoids the reconstruction-error cost paid by a fixed prior.  For neuromorphic hardware under strict power budgets, maximizing CV at the population level may be a more hardware-agnostic objective than enforcing explicit sparsity penalties, whose distributional assumptions need not hold in practice.

\textbf{Limitations and testable predictions.}~\emph{First---the sublinear-cost regime is the framework's most exposed flank.}  The closed-form Gibbs manifold and the ICE--CV monotonicity theorem (Supp.~Sect.~1.4) are proven under \emph{strictly convex} $c(r)$.  Empirical CV$>1$ falls outside the truncated-exponential range $[1/\sqrt{3},1]$ admitted by linear cost on $[0,r_{\max}]$ and is accommodated only by a phenomenological sublinear extension $c(r)\propto r^\alpha$, $\alpha<1$ (Supp.~Note ``Sublinear cost extension''), which violates the convexity premise.  We verify the ICE--CV ordering numerically over the relevant $\alpha$ range, so the high-CV regime in which much of our data sit rests on numerical rather than theorem evidence.  Independent biophysical motivation for an effectively sublinear cost (e.g., shared metabolic infrastructure across active synapses) would strengthen the framework.

\emph{Second}, when noise entropy is non-negligible---as we find on the V1 dataset ($\bar\rho\approx 0.85$)---the Gibbs solution is an upper bound on attainable $I(s;r)/C$ and a coding-readiness measure in spontaneous conditions, not an exact optimum.  Rank ordering and ICE--CV monotonicity within a system are preserved when $\rho$ is approximately uniform across the population (here V1 IQR width $0.06$).

\emph{Third}, the predicted monotonic ICE--CV relationship applies to \emph{population} CV; single-neuron $\mathrm{CV}_L$ can saturate at the refractory period and membrane time constant, as our single-neuron simulations show.

\emph{Fourth}, the V1 data were recorded under sufentanil anesthesia with drifting gratings; whether the manifold structure holds under awake, free-viewing, naturalistic conditions remains to be verified.

\emph{Fifth}, the topological comparison in Fig.~5 uses a per-system $R_0$ rate factor and should be read as trajectory-shape alignment rather than absolute collapse.

\emph{Sixth}, the spike-count agreement in Fig.~4L derives from $P_{\mathrm{agree}}=e^{-2\lambda}I_0(2\lambda)$ for Poisson counts and counts every coincident outcome---including the all-zero ``both silent'' word---as agreement; at very low $\lambda$ its rise with CV is partly silence-induced.  The predicted CV--reliability link is, however, also empirically supported by an FR-controlled decoder on the same V1 data: per-neuron 8-class accuracy correlates positively with $\mathrm{CV}_L$ after partialling FR ($\rho_{\,|\,\mathrm{FR}}=+0.68$, $p=10^{-78}$, pooled $N=564$, 5 sessions; positive in every session, Fisher-combined $\rho=+0.72$, $p=10^{-100}$; Sect.~\ref{sec:network_results}, Supp.~Fig.~S11).  The prediction direction therefore holds beyond the Poisson algebra; decoder-based reliability under awake, naturalistic conditions remains a future test.

\emph{Seventh---high CV is not universally desirable.}  The framework predicts that elevated CV is beneficial \emph{only} under the cost and coding assumptions developed here (non-negative responses, effective sublinear or bounded cost, an information-per-cost coding objective).  In other regimes high variability may degrade fine stimulus discriminability, signal pathological dynamics (e.g., seizure), or reflect noise (degraded sensory periphery, perturbations breaking E--I balance).  The framework therefore identifies a \emph{conditional} normative role for high CV, not a universal ``higher is better'' principle.

The framework generates testable predictions: (i) acutely reducing the energy budget (e.g., cortical cooling or pharmacological suppression of oxidative metabolism) should decrease population CV and sparseness in tandem along the manifold, following the trajectory predicted by the $H$--$C$ curve; (ii) brain regions with tighter metabolic budgets (e.g., higher neuron density relative to vascular supply) should show higher CV and ICE; (iii) developmental sharpening of E--I balance should drive circuits toward the high-CV regime, a trajectory that may reverse with aging or neurodegeneration as inhibitory circuits deteriorate.  Analytical derivations supporting these predictions, including the monotonicity proofs for the ICE--CV and agreement--CV relationships, are provided in Supp.~Sect.~1--2.

\textbf{Conclusion.}~Metabolic constraints actively shape neural response statistics rather than acting only as passive limits.  Within the normative framework developed here, elevated variability is a feature of an energy-efficient operating regime in which information per unit cost is high and sparseness follows from the underlying response distribution; CV serves as the macroscopic coordinate organizing this relationship along the Gibbs manifold.  In biological circuits, E--I balance and synaptic plasticity provide plausible mechanisms for moving activity toward this regime; in artificial systems, directly targeting high CV gives a hardware-agnostic principle for designing energy-efficient representations.
